\section{Tres ejes del sector tecnológico en la bolsa estadounidense: AI, reindustrialización y digitalización financiera}

Después de establecer el contexto temporal, esta sección pasa a la visión de conjunto de Freda sobre los ejes actuales del sector tecnológico en la bolsa estadounidense. Ella considera que hoy hay tres ejes muy claros: AI, Re-industrialization (reindustrialización) y Digitization of Finance (digitalización de la industria financiera). Estos tres ejes no son islas paralelas, sino que están conectados entre sí: la AI necesita centros de datos, energía y chips; la reindustrialización dirige el capital hacia la infraestructura dentro de Estados Unidos; y la digitalización financiera converge con la AI en las monedas estables, los pagos y el Agent Commerce.

\begin{figure}[H]
\centering
\includegraphics[width=0.96\textwidth]{us-tech-three-lines.png}
\caption{Los tres ejes del sector tecnológico en la bolsa estadounidense: la AI, la reindustrialización y la digitalización financiera encajan entre sí. Diagrama conceptual propio, elaborado a partir de la conversación de 00:08:12--00:10:20.}
\end{figure}

\begin{knowledgebox}{Lectura de la figura: los tres ejes no son tres mesas de juego independientes}
La AI necesita chips, nube y centros de datos; la reindustrialización aporta inversión en energía, manufactura e infraestructura; la digitalización financiera entra en la cadena de transacciones mediante las monedas estables, los pagos y el Agent Commerce. Las políticas públicas, los aranceles y las elecciones intermedias determinan la velocidad y la dirección de estos flujos de capital.
\end{knowledgebox}

\subsection{Por qué la AI y la reindustrialización están vinculadas}

Esta sección lleva la AI del relato del software de vuelta a la infraestructura física. Antes se dijo que los tres ejes encajan entre sí; el vínculo más directo son los centros de datos y la inversión en energía: el avance de los modelos requiere capacidad de cómputo, y la capacidad de cómputo requiere electricidad, terrenos, chips y capacidad de construcción.

La capa superficial de la AI son los modelos y las aplicaciones; la capa de fondo son la electricidad, los terrenos, los chips, los centros de datos, las redes y la ingeniería de construcción. Freda menciona que los compromisos de inversión en Estados Unidos de países como Japón y Corea del Sur se destinarán a proyectos de infraestructura como energía y centros de datos, y ese dinero, a su vez, sostendrá la construcción de capacidad de cómputo para AI. Es decir, la AI no es un ciclo puramente de software, sino un ciclo de gasto de capital que arrastra consigo la tecnología dura, la energía, la manufactura y los acuerdos diplomáticos.

\begin{warningbox}{Advertencia de interpretación: las empresas de AI no son SaaS tradicional}
El SaaS tradicional suele tener pocos activos, márgenes brutos altos y costos marginales bajos; las empresas de modelos de AI y parte de las aplicaciones de AI requieren, en cambio, inversión continua en inferencia, entrenamiento, datos e infraestructura. Aplicar directamente a las empresas de AI la plantilla de márgenes brutos del software anterior impide ver el cambio en la estructura de costos.
\end{warningbox}

\subsection{Por qué la digitalización financiera vuelve a ser importante}

A continuación se pasa de la electricidad y los chips al dinero y las transacciones. La digitalización financiera vuelve a ser importante porque, si un Agent realiza compras en nombre del usuario, los pagos, las monedas estables, las autorizaciones y los puntos de entrada a las transacciones pasan a formar parte de la comercialización de la AI.

En este episodio, la digitalización financiera no es la fintech en términos genéricos, sino la combinación de monedas estables, pagos, mercados de predicción, Robinhood y Agent Commerce. Freda menciona que los pagos con monedas estables son cheaper/faster/better, y que en el futuro, cuando un Agent pague y compre cosas en nombre del usuario, los conceptos de token o de moneda estable volverán a plantearse. Este eje conecta los puntos de entrada de la AI, el pago y la liquidación, las transacciones en internet y la regulación financiera.

\begin{importantbox}{Implicaciones financieras del Agent Commerce}
Si el Agent se convierte en el punto de entrada de las compras, no solo afectará a la publicidad y al comercio electrónico, sino también a las billeteras de pago, la autorización de transacciones, las monedas estables, el control de riesgos y la captación de clientes de los comercios. El punto de conexión entre la digitalización financiera y la AI es «quién completa la transacción en nombre del usuario».
\end{importantbox}

\subsection{El lugar del mercado chino a los ojos de los inversionistas estadounidenses}

En el programa, la valoración del mercado chino es bastante mesurada: cada vez más inversionistas estadounidenses visitan China y quedan impresionados por los EV y los robots, pero el monto real de inversión no ha cambiado de forma notable. Lo que más les importa son dos cosas: el resultado de las negociaciones entre China y Estados Unidos, y el dinamismo de las IPO en la bolsa de Hong Kong. Dicho de otro modo, al capital estadounidense no le falta interés, sino que está esperando señales sobre las políticas, las salidas y la calidad de las IPO tecnológicas.

\subsection{Resumen de la sección}

Los tres ejes amplían la AI desde las empresas de modelos hacia la energía, la manufactura, los pagos y las políticas públicas. Esto indica al lector que el ciclo de la AI no es un ciclo único de software, sino un ciclo que redistribuye el gasto de capital, la infraestructura industrial y los puntos de entrada financieros.
